% ==============================================================================
% Modelo para Monografia de Projeto de Graduação (PG)
% Agradecimentos: Prof. Vítor E. Silva Souza - Nemo / DI / UFES
%
% Baseado em abtex2-modelo-trabalho-academico.tex, v-1.9.2 laurocesar
% Copyright 2012-2014 by abnTeX2 group at http://abntex2.googlecode.com/ 
%
% This work may be distributed and/or modified under the conditions of the LaTeX 
% Project Public License, either version 1.3 of this license or (at your option) 
% any later version. The latest version of this license is in
% http://www.latex-project.org/lppl.txt.
%
% IMPORTANTE:
% Instruções encontram-se espalhadas pelo documento. Para facilitar sua leitura,
% tais instruções são precedidas por (*) -- utilize a função localizar do seu
% editor para passar por todas elas.
% ==============================================================================

% Usa o estilo abntex2, configurando detalhes de formatação e hifenização.
\documentclass[
	12pt,				% Tamanho da fonte.
	openright,			% Capítulos começam em página ímpar (insere página vazia caso preciso).
	oneside,			% Para impressão em verso e anverso. Oposto a oneside.
	a4paper,			% Tamanho do papel.
	english,			% Idioma adicional para hifenização.
	french,				% Idioma adicional para hifenização.
	spanish,			% Idioma adicional para hifenização.
	brazil				% O último idioma é o principal do documento.
	]{abntex2}


% Controle das cores.
\usepackage[usenames,dvipsnames]{xcolor}

%%% Importação de pacotes. %%%

\usepackage{forest}

% Conserta o erro "No room for a new \count"
\usepackage{etex}
%\reserveinserts{28}

% Usa a fonte Latin Modern.
\usepackage{lmodern}

% Seleção de códigos de fonte.
\usepackage[T1]{fontenc}

% Codificação do documento em Unicode.
\usepackage[utf8]{inputenc}

% Usado pela ficha catalográfica.
\usepackage{lastpage}

% Indenta o primeiro parágrafo de cada seção.
\usepackage{indentfirst}

% Inclusão de gráficos.
\usepackage{graphicx}

% Inclusão de páginas em PDF diretamente no documento (para uso nos apêndices).
\usepackage{pdfpages}

% Para melhorias de justificação.
\usepackage{microtype}

\usepackage{pstricks,pst-plot,fp}
\usepackage{pstricks-add}
\usepackage{auto-pst-pdf}
\usepackage{subfiles}

% Citações padrão ABNT.
\usepackage[brazilian,hyperpageref]{backref}
\usepackage[alf]{abntex2cite}	
\renewcommand{\backrefpagesname}{Citado na(s) página(s):~}		% Usado sem a opção hyperpageref de backref.
\renewcommand{\backref}{}										% Texto padrão antes do número das páginas.
\renewcommand*{\backrefalt}[4]{									% Define os textos da citação.
	\ifcase #1
		Nenhuma citação no texto.
	\or
		Citado na página #2.
	\else
		Citado #1 vezes nas páginas #2.
	\fi}

% \rm is deprecated and should not be used in a LaTeX2e document
% http://tex.stackexchange.com/questions/151897/always-textrm-never-rm-a-counterexample
\renewcommand{\rm}{\textrm}

% Pacotes não incluídos no template abntex2. 
% Podem ser comentados caso não queira utilizá-los.

% Inclusão de símbolos não padrão.
\usepackage{amssymb}
\usepackage{eurosym}

% Para utilizar \eqref para referenciar equações.
\usepackage{amsmath}

% Permite mostrar figuras muito largas em modo paisagem com \begin{sidewaysfigure} ao invés de \begin{figure}.
\usepackage{rotating}

% Permite customizar listas enumeradas/com marcadores.
\usepackage{enumitem}

% Permite inserir hiperlinks com \url{}.
\usepackage{bigfoot}
\usepackage{hyperref}

% Permite usar o comando \hl{} para evidenciar texto com fundo amarelo. Útil para chamar atenção a itens a fazer.
\usepackage{soulutf8}

\usepackage{algorithm}
\usepackage{algorithmic}

% Colorinlistoftodos package: to insert colored comments so authors can collaborate on the content.
\usepackage[colorinlistoftodos, textwidth=20mm, textsize=footnotesize]{todonotes}
\newcommand{\aluno}[1]{\todo[author=\textbf{Aluno},color=green!30,caption={},inline]{#1}}
\newcommand{\professor}[1]{\todo[author=\textbf{Professor},color=red!30,caption={},inline]{#1}}

% Permite inserir espaço em branco condicional (incluído no texto final só se necessário) em macros.
\usepackage{xspace}

% Permite incluir listagens de código com o comando \lstinputlisting{}.
\usepackage{listings}
\usepackage{caption}
\DeclareCaptionFont{white}{\color{white}}
\DeclareCaptionFormat{listing}{\colorbox{gray}{\parbox{\textwidth}{#1#2#3}}}
\captionsetup[lstlisting]{format=listing,labelfont=white,textfont=white}
\renewcommand{\lstlistingname}{Listagem}
\definecolor{mygray}{rgb}{0.5,0.5,0.5}
\lstset{
	basicstyle=\scriptsize,
	breaklines=true,
	numbers=left,
	numbersep=5pt,
	numberstyle=\tiny\color{mygray}, 
	rulecolor=\color{black},
	showstringspaces=false,
	tabsize=2,
    inputencoding=utf8,
    extendedchars=true,
    literate=%
    {é}{{\'{e}}}1
    {è}{{\`{e}}}1
    {ê}{{\^{e}}}1
    {ë}{{\¨{e}}}1
    {É}{{\'{E}}}1
    {Ê}{{\^{E}}}1
    {û}{{\^{u}}}1
    {ù}{{\`{u}}}1
    {â}{{\^{a}}}1
    {à}{{\`{a}}}1
    {á}{{\'{a}}}1
    {ã}{{\~{a}}}1
    {Á}{{\'{A}}}1
    {Â}{{\^{A}}}1
    {Ã}{{\~{A}}}1
    {ç}{{\c{c}}}1
    {Ç}{{\c{C}}}1
    {õ}{{\~{o}}}1
    {ó}{{\'{o}}}1
    {ô}{{\^{o}}}1
    {Õ}{{\~{O}}}1
    {Ó}{{\'{O}}}1
    {Ô}{{\^{O}}}1
    {î}{{\^{i}}}1
    {Î}{{\^{I}}}1
    {í}{{\'{i}}}1
    {Í}{{\~{Í}}}1
}



%%% Definição de variáveis. %%%

% (*) Substituir os textos abaixo com as informações apropriadas.
\titulo{Classificadores especializados com seleção de atributos SHAP: aplicação no processo Tennessee Eastman e Reator-tanque Agitado Contínuo (CSTR)}
\autor{Gabriel dos Santos Sereno}
\local{Vitória, ES}
\data{2023}
\orientador{Prof. Dr. Thomas Walter Rauber}
\coorientador{Prof. Dr. Celso José Munaro}
\instituicao{
  Universidade Federal do Espírito Santo -- UFES
  \par
  Centro Tecnológico
  \par
  Colegiado do Curso de Engenharia de Computação}
\tipotrabalho{Monografia (PG)}

% Preâmbulo (tipo do trabalho, objetivo, nome da instituição, área de concentração, etc.).
% (*) Verificar se está correto (ex.: substituir por Engenharia de Computação se for o caso).
\preambulo{Dissertação apresentada ao Curso de Inteligência Artificial do Departamento de Informática da Universidade Federal do Espírito Santo, como requisito parcial para obtenção do Grau de Mestre em Inteligência Artificial.}

% Macros específicas do trabalho.
% (*) Inclua aqui termos que são utilizados muitas vezes e que demandam formatação especial.
% Os exemplos abaixo incluem i* (substituindo o asterisco por uma estrela) e Java com TM em superscript.
% Use sempre \xspace para que o LaTeX inclua espaço em branco após a macro somente quando necessário.
\newcommand{\istar}{\textit{i}$^\star$\xspace}
\newcommand{\java}{Java\texttrademark\xspace}
\newcommand{\latex}{\LaTeX\xspace}




%%% Configurações finais de aparência. %%%

% Altera o aspecto da cor azul.
\definecolor{blue}{RGB}{41,5,195}

% Informações do PDF.
\makeatletter
\hypersetup{
	pdftitle={\@title}, 
	pdfauthor={\@author},
	pdfsubject={\imprimirpreambulo},
	pdfcreator={LaTeX with abnTeX2},
	pdfkeywords={abnt}{latex}{abntex}{abntex2}{trabalho acadêmico}, 
	colorlinks=true,				% Colore os links (ao invés de usar caixas).
	linkcolor=blue,					% Cor dos links.
	citecolor=blue,					% Cor dos links na bibliografia.
	filecolor=magenta,				% Cor dos links de arquivo.
	urlcolor=blue,					% Cor das URLs.
	bookmarksdepth=4
}
\makeatother

% Espaçamentos entre linhas e parágrafos.
\setlength{\parindent}{1.3cm}
\setlength{\parskip}{0.2cm}



%%% Páginas iniciais do documento: capa, folha de rosto, ficha, resumo, tabelas, etc. %%%

% Compila o índice.
\makeindex

% Inicia o documento.
\begin{document}


\begin{figure}
  \centering
  \includegraphics[width=.20\textwidth]{figuras/brasao.jpg}
  \label{fig-brasao}
\end{figure}

\begin{center}
  \textbf{\textsf{UNIVERSIDADE FEDERAL DO ESPÍRITO SANTO}}

  \textbf{\textsf{CENTRO TECNOLÓGICO}}

  \textbf{\textsf{COLEGIADO DO CURSO DE ENGENHARIA DE COMPUTAÇÃO}}

  \large{\textbf{\textsf{  }}}

  \large{\textbf{\textsf{  }}}
\end{center}

% Retira espaço extra obsoleto entre as frases.
\frenchspacing

% Capa do trabalho.
\imprimircapa

% Folha de rosto (o * indica que haverá a ficha bibliográfica).
\imprimirfolhaderosto*


% Ficha catalográfica.
% (*) Escolher entre as versões de ficha catalográfica abaixo (comente aquela que não quiser usar).

% Versão 1: caso a biblioteca da sua universidade lhe forneça um PDF (adequar o nome do arquivo).
% \begin{fichacatalografica}
%     \includepdf{include-fichacatalografica.pdf}
% \end{fichacatalografica}

% Versão 2: caso você tenha que inserir sua própria ficha catalográfica.
% (*) Neste caso, preencher palavras-chave e adicione co-orientador (se houver).
\begin{fichacatalografica}
  \vspace*{\fill}
  \hrule
  \begin{center}
    \begin{minipage}[c]{12.5cm}

      \imprimirautor

      \hspace{0.5cm} \imprimirtitulo  / \imprimirautor. --
      \imprimirlocal, \imprimirdata-

      \hspace{0.5cm} \pageref{LastPage} p. : il. (algumas color.) ; 30 cm.\\

      \hspace{0.5cm} \imprimirorientadorRotulo~\imprimirorientador\\

      \hspace{0.5cm}
      \parbox[t]{\textwidth}{\imprimirtipotrabalho~--~\imprimirinstituicao,
        \imprimirdata.}\\

      \hspace{0.5cm}
      1. Diagnóstico de falhas.
      2. Seleção de Características.
      3. Seleção SHAP.
      4. Processo Tennessee Eastman.
      5. Classificadores especializados.
      I. Rauber, Thomas Walter.
      II. Universidade Federal do Espírito Santo.
      IV. \imprimirtitulo \\

      \hspace{8.75cm} CDU 02:141:005.7\\

    \end{minipage}
  \end{center}
  \hrule
\end{fichacatalografica}


% Folha de aprovação.
% (*) Escolher entre as versões de ficha catalográfica abaixo (comente aquela que não quiser usar).

% Versão 1: cópia digitalizada da folha de aprovação assinada pela banca.
% \includepdf{include-folhadeaprovacao.pdf}

% Versão 2: folha de aprovação em branco.
% (*) Ajustar a data e os nomes dos participantes da banca.
\begin{folhadeaprovacao}
  \begin{center}
    {\ABNTEXchapterfont\large\imprimirautor}
    \vspace*{\fill}\vspace*{\fill}
    \begin{center}
      \ABNTEXchapterfont\bfseries\Large\imprimirtitulo
    \end{center}
    \vspace*{\fill}
    \hspace{.45\textwidth}
    \begin{minipage}{.5\textwidth}
      \imprimirpreambulo
    \end{minipage}%
    \vspace*{\fill}
  \end{center}
  Trabalho aprovado. \imprimirlocal, 25 de setembro de 2023:
  \assinatura{\textbf{\imprimirorientador} \\ Orientador}
  \assinatura{\textbf{Professor} \\ Convidado 1}
  \assinatura{\textbf{Professor} \\ Convidado 2}
  %\assinatura{\textbf{Professor} \\ Convidado 3}
  %\assinatura{\textbf{Professor} \\ Convidado 4}
  \begin{center}
    \vspace*{0.5cm}
    {\large\imprimirlocal}
    \par
    {\large\imprimirdata}
    \vspace*{1cm}
  \end{center}
\end{folhadeaprovacao}


% Dedicatória.
% (*) Escrever dedicatória ou remover/comentar seção.
% \begin{dedicatoria}
%    \vspace*{\fill}
%    \centering
%    \noindent
%    \textit{Lorem ipsum dolor sit amet, consectetur adipiscing elit. Duis malesuada laoreet leo at interdum. Nullam neque eros, dignissim sed ipsum sed, sagittis laoreet nisi.} \vspace*{\fill}
% \end{dedicatoria}


% Agradecimentos.
% (*) Escrever agradecimentos ou remover/comentar seção.
\begin{agradecimentos}
  Primeiramente agradecer a Deus que me guiou e fez ser quem eu sou e sou muito grato a tudo o que aconteceu em minha vida. Agradeço também principalmente aos meus pais e ao meu irmão que sempre me apoiaram em tudo e me aconselharam em todas as minhas decisões. A minha esposa Maria que me apoiou e me deu forças para enfrentar todos os tipos de barreiras. Ao o meu professor Dr. Thomas Walter Rauber que me ensinou a ser um futuro mestre e pelo direcionamento para
  desenvolver esse estudo. Ao Dr. Celso José Munaro que também fez parte desse estudo, o meu muito obrigado. Por fim, agradeço a UFES pela oportunidade e também por toda estrutura para o apoio ao estudante.
\end{agradecimentos}


% Epígrafe.
% (*) Escrever epígrafe ou remover/comentar seção.
\begin{epigrafe}
  \vspace*{\fill}
  \begin{flushright}
    \textit{``Quem não vive para servir não serve para viver.''\\
      (Autor Desconhecido)}
  \end{flushright}
\end{epigrafe}


% Resumo em português.
% (*) Escrever resumo e palavras-chave.
\setlength{\absparsep}{18pt}
\begin{resumo}

  Nesse estudo, classificadores especializados são treinados utilizando o método de seleção de característica SHAP no objetivo de identificar falhas nos processos industriais. Nesse contexto, os classificadores recebem dados de entradas a partir da decomposição do problema multi-classe utilizando a técnica One-vs-All. Essa metodologia foi empregada em dois benchmarks muito utilizados no meio academico: Processo Tennessee Eastman e o Reator-tanque Agitado Contínuo. O método SHAP obteve um desempenho de aproximadamente 89\% na métrica F1-score no Processo Tennessee Eastman, selecionando em torno de 4 características em média por falha. Já no Reator-tanque Agitado Contínuo, o algoritmo obteve 90\% na métrica F1-score e selecionou cerca de 3 características em média por falha. Apesar do estudo mostrar que o método Recursive Feature Elimination (RFE) obteve resultados parecidos em comparação ao SHAP em ambos os benchmarks. Entretanto, o método RFE tende a selecionar mais características para ter o mesmo desempenho. Por fim, o estudo sugere que o SHAP pode reduzir a dimensionalidade da base de dados com um bom desempenho na métrica F1-Score, podendo ser utilizado na área industrial e acadêmica.

  \textbf{Palavras-chaves}: Diagnóstico de Falhas. Seleção de Características. SHAP. Processo Tennessee Eastman. Reator-tanque Agitado Contínuo. Classificadores Especializados.
\end{resumo}

% Insere lista de ilustrações.
\pdfbookmark[0]{\listfigurename}{lof}
\listoffigures*
\cleardoublepage

% Insere lista de tabelas.
\pdfbookmark[0]{\listtablename}{lot}
\listoftables*
\cleardoublepage

% Lista de abreviaturas e siglas.
% (*) Preencher com as siglas usadas ao longo do texto e seus significados.
% \begin{siglas}
%   \item[UML] Unified Modeling Language
% \end{siglas}

% Insere o sumário.
\pdfbookmark[0]{\contentsname}{toc}
\tableofcontents*
\cleardoublepage



%%% Início da parte de conteúdo do documento. %%%

% Marca o início dos elementos textuais.
\textual

% Inclusão dos capítulos.
% (*) Para facilitar a organização, os capítulos foram divididos em arquivo separados e colocados dentro da.
% pasta capitulos/. Caso o aluno prefira trabalhar com um só arquivo, basta substituir os comandos \include 
% pelos conteúdos dos arquivos que estão sendo incluídos, excluindo a pasta capitulos/ em seguida.
\subfile{capitulos/ch1-introducao}
\subfile{capitulos/ch2-referencial}
\subfile{capitulos/ch3-requisitos}
\subfile{capitulos/ch4-projeto}
\subfile{capitulos/ch5-conclusoes}

%%% Páginas finais do documento: bibliografia e anexos. %%%

% Finaliza a parte no bookmark do PDF para que se inicie o bookmark na raiz e adiciona espaço de parte no sumário.
\phantompart

% Marca o início dos elementos pós-textuais.
\postextual

% Referências bibliográficas
\bibliography{bibliografia}


% Apêndices.
\begin{apendicesenv}

  % Imprime uma página indicando o início dos apêndices.
  \partapendices

  \begin{table*}[!htbp]
    \centering
    \caption{ Relação entre a falha e as características selecionadas pelo SHAP no CSTR (exceto a característica 13) }
    \begin{tabularx}{\textwidth}{|l|X|X|}
      \hline\textbf{Falha} & \textbf{Descrição da Falha}                                         & \textbf{Características selecionadas}                                                 \\ \hline
      2                    & Obstrução na saída do tanque                                        & Contro. de Nível de Saída \newline Tempe. da E. do Líquido de Refri.                  \\\hline
      3                    & Obstrução no revestimento                                           & Ponto de Ajuste do Líquido de Refri. \newline Pressão da Entrada do Líquido de Refri. \\\hline
      4                    & Vazamento do revestimento para o\newline ambiente                   & Ponto de Ajuste do Líquido de Refri. \newline Vazão do Liquido de Refrigeração        \\\hline
      5                    & Vazamento do revestimento para o\newline tanque                     & Vazão do Liquido de Refrigeração                                                      \\\hline
      6                    & Vazamento da bomba                                                  & Contro. de Nível de Saída \newline Vazão do Produto                                   \\\hline
      7                    & Perda de pressão na bomba                                           & Contro. de Nível de Saída                                                             \\\hline
      8                    & Incrustação na superfície de\newline troca de calor do revestimento & Temperatura do Reator \newline Ponto de Ajuste do Líquido de Refri.                   \\\hline
      10                   & Energia de ativação da reação primária                              & Concentração do Produto A \newline Ponto de Ajuste do Líquido de Refri.               \\\hline
      12                   & Vazão de alimentação anormal                                        & Ponto de Ajuste do Líquido de Refri. \newline Contro. de Nível de Saída               \\\hline
      13                   & Temperatura de alimentação anormal                                  & Temperatura da entrada \newline Ponto de Ajuste do Líquido de Refri.                  \\\hline
      14                   & Concentração de alimentação anormal                                 & Ponto de Ajuste do Líquido de Refri. \newline Vazão do Liquido de Refrigeração        \\\hline
      15                   & Temperatura anormal da água de\newline resfriamento                 & Ponto de Ajuste do Líquido de Refri.                                                  \\\hline
      16                   & Pressão anormal da água de resfriamento                             & Pressão da Entrada do Líquido de Refri. \newline Concentração de entrada              \\\hline
      19                   & Ponto de ajuste anormal do controlador de nível                     & Nível do reator \newline Concentração do Produto A                                    \\\hline
    \end{tabularx}
    \label{tbl-failure-feature-cstr}
  \end{table*}

  \begin{table*}[!htbp]
    \centering
    \caption{ Relação entre a falha e as características selecionadas pelo SHAP no processo Tennessee Eastman }
    \begin{tabularx}{\textwidth}{|l|X|X|}
      \hline
      \textbf{\textbf{Falha}} & \textbf{\textbf{Descrição da Falha}}                  & \textbf{\textbf{Características selecionadas}}                                           \\ \hline
      1                       & Razão de alimentação A/C                              & Alimentação do reagente A (fluxo 1) \newline Vazão de reagente A (fluxo 1)               \\ \hline
      2                       & Composição de B                                       & Taxa de purga (fluxo 9) \newline Válvula de purga (fluxo 9)                              \\ \hline
      3                       & Temperatura de alimentação de D                       & Válvula de vapor do removedor \newline Componente E                                      \\ \hline
      4                       & Temperatura de entrada do refrigerante do reator      & Vazão de refrigerante do reator \newline Temperatura do removedor                        \\ \hline
      5                       & Temperatura de entrada do refrigerante do condensador & Vazão de refrigerante do condensador \newline Fluxo do removedor (fluxo 11)              \\ \hline
      6                       & Perda de alimentação de A                             & Vazão de reagente A (fluxo 1)                                                            \\ \hline
      7                       & Redução na disponibilidade de C                       & Vazão dos reagentes A e C (fluxo 4) \newline Válvula de vapor do removedor               \\ \hline
      8                       & Composição de alimentação de A, B e C                 & Válvula de vapor do removedor \newline Componente E                                      \\ \hline
      9                       & Temperatura de alimentação de D                       & Válvula de vapor do removedor \newline Componente E                                      \\ \hline
      10                      & Temperatura de alimentação de C                       & Vazão de gás do removedor \newline Válvula de vapor do removedor                         \\ \hline
      11                      & Temperatura de entrada do refrigerante do reator      & Vazão de refrigerante do reator \newline Temperatura do reator                           \\ \hline
      12                      & Temperatura de entrada do refrigerante do condensador & Válvula de vapor do removedor \newline Componente E                                      \\ \hline
      13                      & Cinética das reações                                  & Vazão de gás do removedor \newline Válvula de vapor do removedor                         \\ \hline
      14                      & Válvula do Reator da Água de Refrigeração             & Temperatura do reator \newline Temperatura de saída do refrigerante do reator            \\ \hline
      15                      & Válvula do Condensador da Água de Refrigeração        & Temperatura do removedor \newline Vazão de gás do removedor                              \\ \hline
      16                      & Desconhecido                                          & Válvula de vapor do removedor \newline Vazão de gás do removedor                         \\ \hline
      17                      & Desconhecido                                          & Temperatura de saída do refrigerante do reator \newline Temperatura do reator            \\ \hline
      18                      & Desconhecido                                          & Temperatura de saída do refrigerante do separador \newline Válvula de vapor do removedor \\ \hline
      19                      & Desconhecido                                          & Vazão de reciclagem (fluxo 8) \newline Válvula do compressor                             \\ \hline
      20                      & Desconhecido                                          & Válvula do compressor \newline Pressão do separador                                      \\ \hline
      21                      & Válvula (Stream 4)                                    & Vazão dos reagentes A e C (fluxo 4)                                                      \\ \hline
    \end{tabularx}
    \label{tbl-failure-feature-tep}
  \end{table*}

  % (*) Incluir como apêndice a documentação técnica produzida durante o PG (especificação de requisitos,
  % projeto arquitetural, etc.). Utilizar o exemplo \includepdf caso o documento seja produzido em outro
  % editor de texto (Microsoft Word, LibreOffice Writer) e transformado em PDF. Utilizar o exemplo \include
  % caso os documentos tenham sido também escritos em LaTeX.
  % \includepdf[pages={1-}]{apendices/apendice01-requisitos.pdf}
  % \includepdf[pages={1-}]{apendices/apendice02-projeto.pdf}
  % \include{ap1-requisitos}
  % \include{ap2-projeto}
\end{apendicesenv}



% Índice remissivo.
\phantompart
\printindex

% Fim do documento.
\end{document}
